% SEGMENT OLP-0581-S01
\documentclass[../../../include/open-logic-section]{subfiles}

\begin{document}

\olfileid{sth}{cardinals}{cp}
\olsection{காண்டோரின் கொள்கை}

\olref[ordinals][vn]{sec} இல் நாம் பேசியதை நினைவுகூருங்கள். நன்கு
வரிசைப்படுத்தப்பட்ட கணங்களைப் பற்றிப் பேசினோம்; நன்கு வரிசைப்படுத்தல்களின்
பிரதிநிதிகளாக அமையும் பொருள்கள் கிடைத்தால் நன்றாக இருக்கும் என்றும்
கருதினோம். இந்த நோக்கத்துடன் வரிசையெண்களை அறிமுகப்படுத்தினோம். பின்னர்,
\olref[ordinals][ordtype]{ordtypesworklikeyouwant} இல் அவை நாம்
விரும்பியபடியே செயல்படுகின்றன என்று, அதாவது பின்வரும் நிபந்தனையை
நிறைவேற்றுகின்றன என்று, காட்டினோம்:
\[
	\ordtype{A, <} = \ordtype{B, \lessdot}
	\text{ iff } \tuple{A, <} \isomorphic \tuple{B, \lessdot}.
\]

% SEGMENT OLP-0581-S02
இதற்கும் முந்தைய \olref[sfr][siz][equ]{sec} ஐ நினைவுகூருங்கள். அங்கு
கணத்தின் ``அளவு'' என்ற கருத்தை முறைசாரா முறையில் அறிமுகப்படுத்தினோம்.
குறிப்பாக, $f \colon A \to B$ என்ற !!a{bijection} இருக்கும்போதும்
அப்போதுதான் அந்த இரண்டு கணங்களும் எண்ணொத்தவை,
$\cardeq{A}{B}$, என்று சொன்னோம். இந்தக் கருத்து
நன்கு வரிசைப்படுத்தல் என்ற கருத்தைவிட உள்ளார்ந்த வகையில் {எளியது}:
!!{bijection}s இருப்பதில்தான் நமக்கு ஆர்வம்; வரிசை ஓரினச்சார்புகளைப்
போல அந்த !!{bijection}s ``ஏதேனும் கட்டமைப்பைப் பாதுகாக்கின்றனவா''
என்பதில் அல்ல.

% SEGMENT OLP-0581-S03
இவற்றிலிருந்து இயல்பான ஓர் எண்ணம் எழுகிறது. நன்கு வரிசைப்படுத்தல்களை
அளவிடுவதற்காகச் சில பொருள்களாகிய \emph{வரிசையெண்களை} அறிமுகப்படுத்தியதுபோல்,
கணத்தின் அளவை அளவிடுவதற்காகச் சில பொருள்களாகிய \emph{கண அளவெண்களை}
அறிமுகப்படுத்தலாம். இதுவே இந்த அத்தியாயத்தின் நோக்கம்.

கண அளவெண்கள் எவை என்று சொல்வதற்கு முன், அவை நிறைவேற்ற வேண்டிய
கொள்கையை வகுக்க வேண்டும். $X$ என்ற கணத்தின் கண அளவெண்ணை $\card{X}$
என்று எழுதினால், அவை பின்வருமாறு அமைய வேண்டும்:
\[
	\card{A} = \card{B} \text{ iff } \cardeq{A}{B}.
\]
இதை \emph{காண்டோரின்} கொள்கை என்று அழைப்போம்; ஏனெனில் இக்கொள்கையைத்
தெளிவாக மனத்தில் கொண்டிருந்த முதல் நபர் காண்டோராக இருக்கலாம்.
(\olref[hp]{sec} இல் \emph{ஹியூமின்} கொள்கையுடன் இதற்குள்ள தொடர்பைப்
பற்றி மேலும் கூறுவோம்.) ஆகவே, ஒவ்வொரு $X$ க்கும் காண்டோரின் கொள்கை
நிறைவேறும் வகையில் $\card{X}$ ஐ வரையறுப்பதே நமது நோக்கம்.

\end{document}
